\section{Conclusion}
\label{sec:conclusion}

A chain whose block validity requires settlement and DA receipts halts when
those layers fail, although its validators are connected and correct: an
external dependency can cause a liveness failure without a BFT failure.
Engram makes suspending such a dependency an agreed validity and safety mode.
A quorum-backed health certificate drives a four-state FSM above an unchanged
CometBFT agreement core. In the degraded modes the corresponding receipt
requirement is relaxed and cross-chain withdrawals are locked, and the return
to \textsc{Anchored} requires recovery hysteresis, DA backfill of the
\textsc{Sovereign} interval, and a validator-verified recovery proof.

The formal and empirical evidence supports the mechanism within the stated
assumptions. Bounded model checking found no violations in each layer (the
largest run explored \resFSMStatesM{} distinct states), and the cross-layer
refinement remains open. A seven-validator testbed with live Bitcoin and
Celestia networks kept committing blocks under a composite fault and completed
a full recovery cycle with complete AppHash agreement. Against seven vanilla
validators on the same hosts, Engram lengthens the block interval by about
\resOverheadPct{} (\resBlockEngram\,s against \resBlockVanilla\,s).

\subsection*{Artifact Availability}

\pending{repository URL, license, and availability date.}
The implementation, TLA+ specifications, and experiment scripts are
publicly available.
